% =============================================================
%  Gyeongmin Kim — Resume (A4)
%  Build:  tectonic -X compile resume-a4.tex
%  Engine: XeLaTeX (tectonic).  Fonts: xetexko + Noto Sans CJK KR
% =============================================================
\documentclass[a4paper]{article}

% ---- page ---------------------------------------------------
\usepackage[a4paper,top=12mm,bottom=12mm,left=15mm,right=15mm,footskip=5.5mm]{geometry}
\raggedbottom
\setlength{\parindent}{0pt}
\setlength{\parskip}{0pt}

% ---- fonts --------------------------------------------------
\usepackage{fontspec}
\usepackage{xetexko}
% Noto Sans CJK KR for Latin and Hangul alike — one family, so the two
% scripts share weight/height, and the build is portable (no macOS-only
% faces). Install the font locally to build (Ubuntu: fonts-noto-cjk).
% -liga/-rlig kill the f-ligatures: with them on, "Buffer" comes out of
% the PDF as "Buffer" (U+FB00) and ATS keyword matching misses it.
% No italic face exists in the CJK family; FakeSlant covers the odd \emph.
\setmainfont{Noto Sans CJK KR}[
  BoldFont   = Noto Sans CJK KR Bold,
  ItalicFont = *,
  ItalicFeatures = {FakeSlant=0.15},
  Ligatures  = TeX,
  RawFeature = {-liga,-rlig,-clig,-dlig},
]
\setmainhangulfont{Noto Sans CJK KR}[
  BoldFont = Noto Sans CJK KR Bold,
]
\newfontfamily\symfont{Noto Sans CJK KR}

% ---- packages -----------------------------------------------
\usepackage{xcolor}
\usepackage{graphicx}
\usepackage{enumitem}
\usepackage{needspace}
\usepackage{ragged2e}
\usepackage{tikz}
\usepackage{fancyhdr}
\usepackage{lastpage}
\usepackage[hidelinks]{hyperref}

% ---- palette ------------------------------------------------
\definecolor{ink}{HTML}{1A1A1A}
\definecolor{bodyc}{HTML}{333333}
\definecolor{muted}{HTML}{757575}
\definecolor{rulec}{HTML}{C9C9C9}
\definecolor{hairl}{HTML}{E3E3E3}
\definecolor{soft}{HTML}{F5F5F5}
\definecolor{accent}{HTML}{1E3A5F}

% ---- body type ----------------------------------------------
\newcommand{\bodysize}{\fontsize{9}{12.05}\selectfont}
\AtBeginDocument{\color{bodyc}\bodysize}
\RaggedRight
\hbadness=10000
% No automatic hyphenation — it splits identifiers ("FastCon-former").
% \emergencystretch absorbs the odd overlong line instead; breaks at
% hyphens already in the text (gradient-based) stay allowed.
\hyphenpenalty=10000
\exhyphenpenalty=100
\emergencystretch=2em

% ---- helpers ------------------------------------------------
% Arrow glyph, set in \symfont so it stays stable if the main font changes.
% Write `A \ra{} B` for a spaced arrow — the literal spaces stay breakable.
\newcommand{\ra}{{\symfont →}}
\newcommand{\til}{\textasciitilde{}}
\newcommand{\lnk}[2]{\href{#1}{\textcolor{ink}{\bfseries #2}}}
% display type: bold, upright
\newcommand{\disp}{\bfseries}

% ---- footer -------------------------------------------------
\pagestyle{fancy}
\fancyhf{}
\renewcommand{\headrulewidth}{0pt}
\renewcommand{\footrulewidth}{0pt}
\fancyfoot[L]{\color{muted}\fontsize{7.2}{8.6}\selectfont Gyeongmin Kim (김경민)}
\fancyfoot[R]{\color{muted}\fontsize{7.2}{8.6}\selectfont \thepage\,/\,\pageref{LastPage}}

% ---- lists --------------------------------------------------
\newcommand{\dotlabel}{\raisebox{0.5ex}{\tikz\fill[muted] circle (0.5mm);}}
% leftmargin is absolute (lists ignore the outer \leftskip), so it carries
% the full stair depth: entry 5mm -> bullet text 10mm -> sublist 13.4mm.
\setlist[itemize]{
  leftmargin=10mm, labelsep=1.6mm, labelwidth=0pt,
  topsep=0.9mm, itemsep=0.95mm, parsep=0pt, partopsep=0pt,
  label=\dotlabel,
}

% ---- section title ------------------------------------------
\newcommand{\sectiontitle}[1]{%
  \par\leftskip=0pt\relax\addvspace{8.5mm}%
  \needspace{16mm}%
  {\color{ink}\bfseries\fontsize{11.2}{13.4}\selectfont #1}\par
  \vspace{1mm}%
  {\color{rulec}\hrule height 0.4pt}%
  \par\nobreak\vspace{2.4mm}\nobreak\global\justsectiontrue}

% A header must never be the last thing on a page. \needspace inserts a
% \penalty-100, which would let TeX break between a header and its first
% \entry — so headers set this flag and the next \entry binds with
% \nobreak instead, pushing the whole group to the next page if needed.
\newif\ifjusthead
% Set by \sectiontitle so the first \company under it skips its separator
% rule — the section's own accent rule already does that job there.
\newif\ifjustsection

% ---- company entry ------------------------------------------
% \company{name}{dates}{affiliation}
% A separator rule above the name marks "new company starts here" — pure
% whitespace would need ~5mm per block to read as a break, which costs a page.
\newcommand{\company}[3]{%
  \par\leftskip=0pt\relax
  \ifjustsection
    \nobreak
  \else
    \addvspace{3.2mm}\needspace{18mm}%
    {\color{rulec}\hrule height 0.5pt}\nobreak\vspace{2.6mm}%
  \fi
  \global\justsectionfalse
  \noindent\hbox to \linewidth{%
    {\color{ink}\disp\fontsize{10.2}{12.2}\selectfont #1}%
    \hfill
    {\color{muted}\bfseries\fontsize{8.6}{10.3}\selectfont #2}}\par
  \vspace{0.5mm}%
  {\color{muted}\fontsize{8.4}{10.5}\selectfont #3}\par
  \vspace{1.3mm}%
  {\color{rulec}\hrule height 0.5pt}%
  \par\nobreak\vspace{2.2mm}\nobreak\global\justheadtrue}

% ---- portfolio category header ------------------------------
\newcommand{\catheader}[1]{%
  \par\leftskip=0pt\relax
  \ifjustsection\nobreak\else\addvspace{2.6mm}\needspace{12mm}\fi
  \global\justsectionfalse
  {\color{ink}\disp\fontsize{10.2}{12.2}\selectfont #1}\par
  \vspace{1.1mm}%
  {\color{rulec}\hrule height 0.5pt}%
  \par\nobreak\vspace{2.2mm}\nobreak\global\justheadtrue}

% ---- project entry inside a company -------------------------
\newcommand{\entry}[1]{%
  \par\leftskip=0pt\relax
  \ifjusthead\nobreak\else\ifjustsection\nobreak\else\addvspace{3.4mm}\needspace{9mm}\fi\fi
  \global\justheadfalse\global\justsectionfalse
  % Headings all sit flush at the margin — level is carried by type size
  % (11.2 section > 10.2 company > 9.7 entry). Only the entry's elements
  % (description, keywords; bullets via list leftmargin) indent.
  \noindent\raisebox{0.62ex}{\tikz\fill[muted] (0,0) rectangle (1.3mm,1.3mm);}\hspace{1.5mm}%
  {\color{ink}\disp\fontsize{9.7}{12.2}\selectfont #1}%
  \par\nobreak\vspace{1.4mm}\nobreak
  \leftskip=5mm\relax}

% Sub-level for a project's implementation details. Indentation is only
% worth its horizontal cost when the levels are genuinely different —
% here: entry (what area) > bullet (which project) > dash (how it works).
\newcommand{\subdash}{\raisebox{0.44ex}{\color{muted}\rule{1.5mm}{0.4pt}}}
\newlist{sublist}{itemize}{1}
\setlist[sublist]{
  leftmargin=3.4mm, labelsep=1.5mm, labelwidth=0pt,
  topsep=0.5mm, itemsep=0.5mm, parsep=0pt, partopsep=0pt,
  label=\subdash,
}

% OVERVIEW row: fixed-width label column, wrapping body that stays aligned
% under itself. Deliberately an index, not half-sentences with a dash.
\newcommand{\ovrow}[2]{%
  \noindent\hbox to 11mm{{\color{ink}\bfseries #1}\hfil}%
  \parbox[t]{\dimexpr\linewidth-11mm\relax}{\RaggedRight #2}\par}

\newcommand{\ovrowc}[2]{%
  \noindent\hbox to 13mm{{\color{ink}\bfseries #1}\hfil}%
  \parbox[t]{\dimexpr\linewidth-13mm\relax}{\RaggedRight #2}\par}

% Team size / role rides on the tail of the description paragraph rather
% than owning a bullet — the last line of a paragraph usually has room,
% so this costs no extra lines.
\newcommand{\team}[1]{\ {\color{muted}(#1)}}

% ---- publications -------------------------------------------
% Unaccepted work must not sit directly under the "Publications" heading —
% every source consulted agrees (Philosophers' Cocoon: "under no
% circumstances"; Academia SE: "creates an impression of ... exaggeration").
% Hence the Preprints / Under review sub-labels.
\newif\ifjustpubgroup
\newcommand{\pubgroup}[1]{%
  \par\leftskip=0pt\relax\addvspace{2.6mm}%
  \needspace{15mm}%
  {\color{muted}\bfseries\fontsize{8}{9.6}\selectfont #1}\par
  \nobreak\vspace{1.1mm}\nobreak\global\justpubgrouptrue}

% \pubentry{title}{authors + status + links}
% Citation only, no description. Of 11 real industry CVs measured, 8 carry
% no description at all and the other 3 carry exactly one line. The impact
% lives in Experience/Projects instead (Rolland, Jiao, Asyrofi all do this).
\newcommand{\pubentry}[2]{%
  \par
  \ifjustpubgroup\nobreak\else\addvspace{2mm}\needspace{13mm}\fi
  \global\justpubgroupfalse
  {\color{ink}\disp\fontsize{9.7}{12.2}\selectfont #1}\par
  \nobreak\vspace{0.7mm}%
  {\fontsize{8.6}{10.9}\selectfont #2}\par}

% ---- keyword strip ------------------------------------------
% \parbox runs \@parboxrestore, which resets \rightskip — so re-assert
% \RaggedRight inside, or the strip comes out justified and stretches.
\newcommand{\keywords}[1]{%
  \par\vspace{1.5mm}%
  {\color{muted}\fontsize{7.4}{9.4}\selectfont #1\par}%
  \par\addvspace{0.3mm}}
\setlength{\fboxsep}{1.6mm}

% =============================================================
\begin{document}

% ------------------------------- HERO ------------------------
\noindent
{\color{ink}\disp\fontsize{20}{21}\selectfont Gyeongmin Kim (김경민)}\par
\vspace{2.3mm}
% One hero paragraph, no keyword table — the whole resume compressed to
% identity -> data/model/serving arc with endpoints -> proof.
{\fontsize{9.2}{12.8}\selectfont TTS · ASR · VC · SE · LLM 모델을 데이터 구축부터 모델 최적화·추론 가속, 서빙까지 직접 다룹니다.}\par
\vspace{3mm}
{\color{hairl}\hrule height 0.5pt}
\vspace{2.2mm}
{\fontsize{7.6}{9.6}\selectfont
\lnk{mailto:kdrkdrkdr@hanyang.ac.kr}{kdrkdrkdr@hanyang.ac.kr} \,·\, \textbf{(+82) 10-3323-0349} \,·\, \lnk{https://github.com/kdrkdrkdr}{github.com/kdrkdrkdr} \,·\, \lnk{https://www.linkedin.com/in/kdrkdrkdr}{linkedin.com/in/kdrkdrkdr} \,·\, \textbf{병역: 산업기능요원 편입 희망}}\par

\par\vspace{4mm}
{\color{ink}\hrule height 1.4pt}

% --------------------------- PUBLICATIONS --------------------
\needspace{34mm}%
\sectiontitle{Publications}

\pubgroup{Preprints}

\pubentry{faster-enhancer.c: A Dependency-Free int8 Runtime for Streaming Speech Enhancement on Commodity CPUs}
{\textbf{Gyeongmin Kim}. arXiv:2607.25350, 2026.07. \enspace [\,\lnk{https://arxiv.org/abs/2607.25350}{arXiv}\,] [\,\lnk{https://github.com/kdrkdrkdr/faster-enhancer.c}{code}\,]}

\pubentry{Extracting Voice Styles from Frozen TTS Models via Gradient-Based Inverse Optimization}
{\textbf{Gyeongmin Kim}. arXiv:2607.25351, 2026.07. \enspace [\,\lnk{https://arxiv.org/abs/2607.25351}{arXiv}\,] [\,\lnk{https://github.com/kdrkdrkdr/supertonic.embed}{code}\,]}

\pubgroup{Under review}

\pubentry{Why One Small Fixed Tree Suffices: The Cost Geometry of Draft Trees for a Multi-Token-Prediction Drafter}
{\textbf{Gyeongmin Kim}, Ayoung Moon, Seung Jin Lee. Under review, 2026.}

% ------------------------------ C.V. -------------------------
\sectiontitle{Experience}

\company{\lnk{https://nc-ai.com}{NC AI} (KR)}{2026년 5월 \til{} 현재}{소속: AI데이터실 데이터엔지니어링팀 (팀원)}

\entry{[콘텐츠진흥원 과제] 문화 콘텐츠 추천 대화 시스템 설계 및 VAETKI 양자화·서빙}
외국인·문화 소외계층에게 한국 문화 콘텐츠를 추천하는 대화 시스템의 LLM 파이프라인을 설계·구현하고, VAETKI-120B-A10B의 양자화와 서빙을 담당했습니다.\team{프로젝트 인원: \mbox{6명} · 역할: 시스템 설계·개발, 모델 양자화, 서빙}\par
\begin{itemize}
  \item \textbf{이원 경로 설계} --- 긴 프롬프트에서 지시 이행이 무너지는 모델 한계를 경로 분리로 우회 --- 일반 대화와 추천(전용 프롬프트 + RAG API 연동) 요청을 별도 VAETKI 호출로 나눠 경로당 프롬프트를 최소화
  \item \textbf{컨트롤러} --- Qwen3.5-4B가 실시간 대화에서 페르소나(취향·호불호)를 추출하고 사용자의 한국어 수준(상/하)을 분류해 VAETKI 응답에 반영, 대화 맥락으로 추천 시점을 판단해 추천 경로로 전환
  \item \textbf{통합 서버} --- 전체 흐름을 FastAPI 서버로 구현, 응답을 문장 단위로 분할해 TTS API로 비동기 스트리밍 전달
  \item \textbf{양자화} --- attention · FFN · MoE expert의 전 Linear를 INT8로 내리고, 손실 민감 지점(lm\_head · router gate · RMSNorm · embedding)은 BF16 가중치로 유지(추론 시 fp32 연산) --- \textbf{상대 L2 오차 0.093\% · 코사인 유사도 0.99961 · SNR 40.9dB}
  \item \textbf{서빙} --- 양자화 모델을 vLLM으로 서빙, Demo 페이지 구축
\end{itemize}
\keywords{키워드: Python, LLM, RAG, Quantization, W8A16, INT8, MoE, vLLM, FastAPI, Streaming, Model-Serving, Inference-Optimization}

\entry{[데이터 생성 파이프라인] 구축 및 추론 가속}
대량의 데이터를 생성하는 파이프라인을 구축하고, 추론 처리량 병목을 \textbf{노드당 처리량}과 \textbf{노드 간 분산} 두 축으로 해결했습니다.\team{프로젝트 인원: \mbox{3명} · 역할: 추론 최적화 · 연구 리드 — 결과 논문 투고}\par
\begin{itemize}
  \item \textbf{노드당} --- Gemma4-31B용 multi-token-prediction assistant 모델을 drafter로 활용해 speculative decoding을 사내 파이프라인에 도입, vLLM · A100 x4 환경에서 \textbf{AR 대비 2\til4$\times$} 가속으로 운영 중
  \item \textbf{노드 간} --- 구성원들이 각자 띄운 vLLM 서버는 중계 계층이 없어 매번 IP·포트로 수동 라우팅해야 했고, 자동 분산 게이트웨이를 제안·구현해 GPU 4장 서버 8대(총 32 GPU)를 주소 등록만으로 단일 엔드포인트에 통합
  \begin{sublist}
    \item 게이트웨이 자체가 병목이 되지 않도록 Go로 구현 --- \textbf{동시 요청 1,000건 기준 라우팅 지연 평균 15ms}. OpenAI API 호환 인터페이스로 기존 클라이언트 코드 수정 없이 전환, 현재 사내 운영 중
  \end{sublist}
\end{itemize}
\keywords{키워드: Python, Go, LLM, Speculative-Decoding, vLLM, API-Gateway, Load-Balancing, Concurrency, OpenAI-Compatible-API, Data-Generation, Inference-Optimization}

\company{\lnk{https://dr-you-group.github.io}{연세대학교 의료원 산학협력단} (KR)}{2025년 3월 \til{} 10월}{소속: 연세대학교 의생명시스템정보학교실 유승찬 교수 연구실 (연구원)}

\entry{[NGS 기반 임상리포트 자동생성 파일럿 사업 (국책과제)] 개발 리드}
엑셀 파일의 검사결과를 NGS 결과지 템플릿에 수작업으로 옮기는 과정에서 휴먼 에러 가능성이 커, 이를 자동화하는 웹 기반 리포트 생성 도구를 개발했습니다.\team{프로젝트 인원: \mbox{3명} (의생명시스템정보학교실/병리과) · 역할: 개발, QA}\enspace[\,\lnk{https://github.com/dr-you-group/NGS-E2E-Pipeline}{code}\,]\par
\begin{itemize}
  \item 엑셀 파일 업로드 시 Web PDF로 자동 변환·저장하는 기능 추가
  \item 각 컴포넌트의 요소, 유전체 검사 결과 값 및 위치, 코멘트 수정 기능 추가
  \item SQLite+Json+Javascript로 데이터 파싱 및 페이지 분할 기능 구현
\end{itemize}
\keywords{키워드: Python, FastAPI, Medical-Data-Processing, Html/CSS/Javascript}

\entry{[SICU False Alarm Monitoring System] 개발 리드}
외과계 중환자실의 잦은 알람(가짜 알람 포함)으로 간호사들이 극심한 스트레스를 받고 있어, 알람의 진위를 판별해 줄이려는 연구에 쓰일 데이터 뷰어 프로그램을 개발했습니다.\team{프로젝트 인원: \mbox{5명} (의생명시스템정보학교실/마취통증의학과) · 역할: 개발, QA}\enspace[\,\lnk{https://github.com/dr-you-group/SICU-Alarm-Monitoring}{code}\,]\par
\begin{itemize}
  \item 입원 시간별 날짜 선택(입원기간/날짜)
  \item 환자 위태 심각도에 따른 색깔별 알람 선택 기능
  \item 간호 기록 엑셀화 및 필터링 기능 (전후 30분 간호중재/간호활동/간호속성코드/속성/Duty/시행일시)
  \item 알람 울린 시간대의 환자 파형 뷰 (ABP/Lead-II/Resp/Pleth)
  \item 알람 진위 판정, 코멘트 작성·저장 기능
\end{itemize}
\keywords{키워드: Python, Medical-Data-Processing, PySide6, SQLite}

\company{\lnk{https://www.ncsoft.com}{엔씨소프트} (KR)}{2024년 7월 \til{} 2025년 1월}{소속: AI데이터실 오디오데이터팀 (팀원)}

\entry{[AUDT\_자동화 프로젝트] 오디오 후처리 자동화 프로그램 개발}
상업용 오디오 처리 프로그램이 원본 파형을 손상시켜 팀원들이 데이터를 수작업으로 처리하고 있어, 품질을 유지하며 후처리를 자동화하는 프로그램을 개발했습니다.\team{프로젝트 인원: \mbox{5명} · 역할: 개발, QA}\par
\begin{itemize}
  \item 녹음실 미세소음 제거: 녹음된 성우의 음성 파일에서 녹음실의 미세한 잡음을 남김없이 제거
  \item 숨소리 제거: F0 미검출 특성과 보정값을 이용하여 음성에서 숨소리를 제거
  \item Endpoint Detection: TTS 학습을 위해 맨 앞과 뒷부분을 감지하여 공백을 일정하게 삽입
  \item 텍스트 태깅: 대만어 데이터의 간체자 혼입을 태깅하여 대사 수정을 빠르게 할 수 있도록 구현
  \item 문장 부호 공백 삽입: CTC forced aligner로 오디오-텍스트를 단어 단위 정렬해, 문장 부호 종류별로 다른 길이의 공백을 삽입
\end{itemize}
\keywords{키워드: Text-to-Speech, Data Processing Automation, CTC Forced-Alignment, PyTorch, GUI, Numpy}

\company{\lnk{https://taiyakistudios.com}{Taiyaki Studios Ltd.} (US)}{2023년 1월 \til{} 7월}{소속: Artificial Intelligence (팀원)}

\entry{[사내용 TTS 학습 도구 개발]}
사내 구성원의 목소리로 TTS를 만들려 했으나 음성 데이터가 적고 파이프라인도 없어, 적은 데이터로도 학습 가능한 TTS 학습 도구를 개발했습니다.\team{프로젝트 인원: \mbox{2명} · 역할: 개발, QA}\par
\begin{itemize}
  \item 데이터셋 구축: 긴 음성 파일을 TTS 학습에 적합하게 분할하고 대사 파일을 자동 생성. 음성 변환 및 증강 기법으로 데이터 부족 문제 해결
  \item 모델 학습: TTS 모델과 G2P 코드를 개발하고, 사전 학습 모델 기반으로 적은 데이터로도 미세 조정이 가능하도록 구현. 모델 압축을 통해 저장 효율성 향상
  \item 성능 테스트: RTF로 추론 속도를, MOS로 발음 정확성을 측정. 하이퍼파라미터 조정으로 모델을 경량화하여 빠른 추론 속도 달성
  \item 파이프라인 구축: TTS 모델을 실제 서비스에 배포하기 위한 추론 코드를 작성하여 파이프라인 구축
\end{itemize}
\keywords{키워드: Python, Text-to-Speech, Few-shot, Voice-Conversion, Deep-Learning}

% --------------------- PERSONAL PROJECT ----------------------
\sectiontitle{Personal Projects}

\entry{모델 최적화 (2024.11 \til{} 현재)}
PyTorch/ONNX 의존성 없이 음성 모델을 순수 C로 재구현하고 최적화하여, GPU/NPU 없이 CPU만으로 실시간 동작하게 만들었습니다.\par
\begin{itemize}
  \item \lnk{https://github.com/kdrkdrkdr/nemotron-asr-streaming.c}{nemotron-asr-streaming.c} --- 0.6B 스트리밍 ASR (FastConformer + RNN-T) [\,\lnk{https://github.com/kdrkdrkdr/nemotron-asr-streaming.c/blob/w8a8/MODEL.md}{구현 문서}\,] [\,\lnk{https://huggingface.co/kdrkdrkdr/nemotron-3.5-asr-streaming-0.6b-w8a8}{양자화 모델}\,]
  \begin{sublist}
    \item cache-aware 스트리밍 상태 유지, chunk 지연 5단계(80\til1120ms, 기본 320ms) 선택
    \item W8A8 양자화 --- per-row int8 + 32-bit row scale를 4-row 타일로 packing(``Q8P''), 가중치 0.62GiB
    \item NEON/AVX2 SIMD 자동 디스패치, 최대 16 스레드 병렬화
  \end{sublist}
  \item \lnk{https://github.com/kdrkdrkdr/lilac}{LILAC} --- 3초 레퍼런스만으로 동작하는 zero-shot 실시간 음성 변환 (OpenVoice v2 포팅) [\,\lnk{https://github.com/kdrkdrkdr/lilac/blob/main/Optimization.md}{최적화 리스트}\,]
  \begin{sublist}
    \item 스트리밍 HiFi-GAN 디코더 --- per-layer state 유지로 hop당 신규 프레임만 생성
    \item 2-thread SPSC ring 오디오 파이프라인, 3-way 병렬 resblock dispatch, RNNoise SIMD(AVX2) 통합
    \item AGC + VAD hangover gate. 16-core CPU에서 RTF 0.7\til0.8, end-to-end 270\til330ms
  \end{sublist}
  \item \lnk{https://github.com/kdrkdrkdr/MossTTS-Nano.c}{MossTTS-Nano.c} --- 100M TTS 모델 [\,\lnk{https://github.com/kdrkdrkdr/MossTTS-Nano.c\#optimizations-applied}{최적화 리스트}\,]
  \begin{sublist}
    \item NEON/AVX SIMD · KV cache · pthread 병렬화 등 12단계 최적화로 합성 68s \ra{} 2.3s
    \item PyTorch CPU 대비 1.8배 빠른 RTF 0.33
  \end{sublist}
  \item \lnk{https://github.com/kdrkdrkdr/DeepFilterNet3.c.wasm}{DeepFilterNet3.c.wasm} --- 노이즈 제거 모델의 순수 C 재작성 + WASM 빌드 [\,\lnk{https://github.com/kdrkdrkdr/DeepFilterNet3.c.wasm/blob/main/OPTIMIZATION.md}{최적화 리스트}\,]
  \begin{sublist}
    \item 기존 WASM 빌드가 모바일에서 실시간 처리가 안 되던 것을 브라우저에서 동작하도록 최적화 --- 프레임당 MacBook M2 \til1ms · Galaxy S23 \til4ms
  \end{sublist}
\end{itemize}
\keywords{키워드: C, WebAssembly, SIMD, OpenBLAS, Quantization, Streaming-ASR, Voice-Conversion, Noise-Reduction, Real-time-Audio-Processing, Python}

\entry{\lnk{https://www.axcellworks.com}{Axcellworks} (JP) 실시간 음성 변환 명료도 개선 --- 원격 기술 협업 (2023.11 \til{} 2024.02)}
일본 스타트업 Axcellworks와 원격으로 협업해, 실시간 음성 변환에서 청크 단위로 추론할 때 경계마다 목소리가 끊기고 명료도가 떨어지는 문제를 개선했습니다.\par
\begin{itemize}
  \item 청크 경계에서 모델의 receptive field가 잘려 컨텍스트가 부족해지는 것이 원인 --- 좌우 이웃 청크를 컨텍스트로 덧대 추론하고 \textbf{가운데 청크만 취하는 overlap-save 방식}으로 경계 아티팩트를 제거, 기존 대비 지연 증가 없이 변환 음성의 명료도 개선
\end{itemize}
\keywords{키워드: Realtime-Voice-Conversion, Python}

\entry{다국어 음성 합성 (2022.10 \til{} 2025.02)}
학습 데이터가 부족한 조건에서 VITS를 파생시켜 언어와 화자를 늘리고, 캐릭터 TTS 서비스로 배포했습니다.\par
\begin{itemize}
  \item \textbf{모델} --- VITS 파생 3종
  \begin{sublist}
    \item 일본어 데이터셋만으로 19개국어 발화 유도 (\lnk{https://github.com/kdrkdrkdr/JA2ML-VITS}{JA2ML-VITS})
    \item 한국어/일본어 이중 언어 (\lnk{https://github.com/kdrkdrkdr/JK-VITS}{JK-VITS})
    \item RVC로 음성 데이터셋을 변환해 학습 데이터 확보 (\lnk{https://github.com/kdrkdrkdr/RVC-VITS}{RVC-VITS})
  \end{sublist}
  \item \textbf{G2P} --- pyopenjtalk 내부 API 수정 + g2pk3 + hangulize를 조합한 음소 변환 체인으로 JA2ML-VITS의 다국어 발화를 구현, 자작 패키지 2종은 PyPI로 공개
  \begin{sublist}
    \item \lnk{https://github.com/kdrkdrkdr/g2pk3}{g2pk3} --- 한국어/일본어/영어를 한국 발음으로 표기 [\,\lnk{https://pypi.org/project/g2pk3/}{PyPI}\,]
    \item \lnk{https://github.com/kdrkdrkdr/ko2kana}{ko2kana} --- 한국어/영어 발음을 가타카나로 변환 [\,\lnk{https://pypi.org/project/ko2kana/}{PyPI}\,]
  \end{sublist}
  \item \textbf{서비스} --- HuggingFace Space 2종
  \begin{sublist}
    \item 학습 데이터 --- BGM 제거(Kim Vocal) \ra{} 화자 분리(pyannote) \ra{} VAD 분할 \ra{} 전사(Whisper)로 3\til10초 클립 약 3만 개 구축
    \item \lnk{https://huggingface.co/spaces/kdrkdrkdr/ProsekaTTS}{ProsekaTTS} --- 2026년 07월 기준 \textbf{누적 방문 230만 · 개발자 복제(Duplicate) 60회}
    \item \lnk{https://huggingface.co/spaces/kdrkdrkdr/ShirokoTTS}{ShirokoTTS} --- 영어/한국어/일본어 3개국어 지원
  \end{sublist}
\end{itemize}
\keywords{키워드: Python, Text-to-Speech, Fine-Tuning, Grapheme-to-Phoneme(G2P), Phoneme-Conversion, Voice-Conversion, Source-Separation, Speaker-Diarization, Whisper, Data-Augmentation, Selenium, Distributed-Crawling}

\entry{일본어 번역 도구 (2020.12 \til{} 2022.04)}
일본어 콘텐츠를 한국어로 번역하는 다양한 도구를 개발했습니다.\par
\begin{itemize}
  \item \lnk{https://github.com/kdrkdrkdr/novel_reader}{novel-reader} --- 소설가가되자, 카쿠요무 등 7개 사이트의 소설을 다국어로 번역하는 안드로이드 앱. 고유명사 사전 등록 시스템 탑재
  \item \lnk{https://github.com/kdrkdrkdr/EhndWebTranslate}{EhndWebTranslate} --- 일본어 웹사이트를 비동기 번역하여 표시. 문서/소설/실시간 번역 지원
  \item \lnk{https://github.com/kdrkdrkdr/UserDict4Papago}{UserDict4Papago} --- Papago 한일 번역의 고유명사 문제를 사전 등록으로 개선
\end{itemize}
\keywords{키워드: Android, Java, Python, Machine-Translation, PySide2, Flask, Selenium}

% ------------------- EXTRACURRICULAR -------------------------
\sectiontitle{Extracurricular Activities}

\company{\lnk{https://elnino.kr}{엘니뇨} (KR)}{2024년 4월 \til{} 2026년 4월}{역할: 대표 (개발, QA, 서류 작성 · 2인 팀) · 창업동아리 시작(2024.04) \ra{} 개인사업자 등록(2025.02) \ra{} 운영 종료}

\entry{[Knoc] 실시간 번역 자막 서비스}
문맥과 분야를 인식해 정확하고 빠르게 번역하는 수업·컨퍼런스용(온/오프라인) 통역기를 개발했습니다.\par
\begin{itemize}
  \item Python 기반 초기 서비스를 대규모 웹소켓 처리를 위해 Go로 전면 재작성
  \item Redis Pub/Sub 기반 실시간 번역 방 및 스크린 송출 시스템 구축
  \item 분산 컴퓨팅 아키텍처 설계 및 구현
  \item 웹 클라이언트 개발: 유튜브형/컨퍼런스형 자막 모드, Context \& Field-Based Any-to-Any Translation Model 적용
  \item 사용자 대시보드 개발: 사용시간/통역시간 관리, PAYG 및 구독제 결제·플랜·영수증 관리, 팀 멤버 관리, 사용 기록 조회
  \item 관리자 대시보드 개발: 팀 정보 관리, 견적서 입력, 서버 리소스 관리 및 로그 모니터링, 팀별 플랜 생성·결제 관리, 팀 통역시간 실시간 모니터링
\end{itemize}
\keywords{키워드: Go, Speech-to-Text, LLM, Prompt-Engineering, Redis, PostgreSQL, NginX, WebSocket, Distributed-Computing, Html/CSS/Javascript}

\entry{[Knoc-Overlay] 데스크톱 자막 오버레이 애플리케이션}
Knoc 실시간 번역 자막을 데스크톱 화면 위에 오버레이로 표시하는 Electron 기반 애플리케이션을 개발했습니다.\par
\keywords{키워드: Electron, Javascript, Html/CSS}

\entry{\lnk{https://elnino.kr}{elnino.kr} 회사 웹사이트 개발}
엘니뇨 공식 웹사이트를 직접 디자인 및 개발했습니다.\par
\keywords{키워드: Html/CSS/Javascript}

\entry{관련 활동}
\begin{itemize}
  \item \textbf{\lnk{https://multilingual.com/elnino-launches-knoc-ai-subtitle-solution-multilingual-communication/}{Multilingual.com 해외 언론 보도 (2026.04)}}
  \item \textbf{실시간 번역 자막 서비스 Knoc 정식 출시 (2026.03)}
  \item 건국대학교 윈터스쿨 번역 자막 서비스 제공 (2026.02)
  \item 대한영상의학기술학회 실시간 번역 자막 서비스 제공 (2025.03)
  \item \textbf{엘니뇨 창업 (2025.02)}
  \item \textbf{교육부 U300+ 학생 창업 유망팀 성장트랙 최종 선발 (2024.08)}
  \item 한양대학교 청년 창업아이템 챌린지 경진대회 본선 발표 (2024.06)
  \item 한양대학교 창업동아리 경진대회 본선 발표 (2024.06)
  \item Naver D2SF 서류합격 (2024.05)
  \item 한양대학교 창업동아리 엘니뇨 대표 (2024.04)
\end{itemize}

% ------------------ LICENSE / EDUCATION ----------------------
\sectiontitle{License \& Education}

\noindent
\begin{minipage}[t]{\dimexpr0.5\linewidth-3mm\relax}\vspace{0pt}\RaggedRight%
  {\color{ink}\disp\fontsize{10.2}{12.2}\selectfont License}\par
  \vspace{1.4mm}
  \begin{itemize}[itemsep=0.9mm,topsep=0pt,leftmargin=3.2mm]
    \item 프로그래밍기능사 (구 정보처리기능사) (2026.07.03)
    \item 운전면허 2종 (2023.07.31)
    \item ITQ A등급 (2016.12.15)
  \end{itemize}
\end{minipage}\hspace{6mm}%
\begin{minipage}[t]{\dimexpr0.5\linewidth-3mm\relax}\vspace{0pt}\RaggedRight%
  {\color{ink}\disp\fontsize{10.2}{12.2}\selectfont Education}\par
  \vspace{1.4mm}
  \begin{itemize}[itemsep=0.9mm,topsep=0pt,leftmargin=3.2mm]
    \item \lnk{http://cs.hanyang.ac.kr/}{한양대학교 공과대학 컴퓨터소프트웨어학부}\\
          {\color{muted}\fontsize{8.1}{10}\selectfont 재학: 2023.03 \til{} 2024.06, 휴학: 2024.07 \til{}}
  \end{itemize}
\end{minipage}

\end{document}
